% Part: incompleteness
% Chapter: representability-in-q
% Section: undecidability

\documentclass[../../../include/open-logic-section]{subfiles}

\begin{document}

\olfileid[hi]{inc}{req}{und}
\olsection{अनिर्णेयता}

हम किसी सिद्धांत $\Th{T}$ को \emph{अनिर्णेय} कहते हैं यदि ऐसी कोई
संगणनात्मक प्रक्रिया न हो जो~$\Th{T}$ की भाषा के किसी भी वाक्य~$!A$ के लिए
परिमित चरणों के बाद, बिना विफल हुए, इस प्रश्न का सही उत्तर दे: ``क्या
$\Th{T}$,~$!A$ को सिद्ध करता है?'' अतः $\Th{Q}$ तभी और केवल तभी निर्णेय
होता जब ऐसी कोई संगणनात्मक प्रक्रिया होती जो अंकगणित की भाषा का कोई
वाक्य~$!A$ दिए जाने पर तय करती कि $\Th{Q} \Proves !A$ है या नहीं। हम इसे
यह पूछकर अधिक सटीक बना सकते हैं: क्या संबंध~$\Prov[\Th{Q}](y)$, जो~$y$
के लिए तभी लागू होता है जब $y$,~$\Th{Q}$ में सिद्ध किए जा सकने वाले किसी
वाक्य की ग्योडेल (G\"odel) संख्या हो, पुनरावर्ती है? उत्तर है: नहीं।

\begin{thm}
$\Th{Q}$ अनिर्णेय है, अर्थात् संबंध
\[
\Prov[\Th{Q}](y) \defiff \fn{Sent}(y) \land
\lexists[x][\Prf[\Th{Q}](x, y)]
\]
पुनरावर्ती नहीं है।
\end{thm}

\begin{proof}
मान लें कि वह पुनरावर्ती है। तब हम विराम समस्या को इस प्रकार हल कर सकते
हैं: $e$ और $n$ दिए जाने पर हम जानते हैं कि
$\cfind{e}(n) \fdefined$ तभी और केवल तभी जब कोई~$s$ ऐसा हो कि
$T(e, n, s)$, जहाँ $T$,~\olref[cmp][rec][nft]{thm:kleene-nf} का क्लेनी
(Kleene) विधेय है। चूँकि $T$ आदिम पुनरावर्ती है, उसका निरूपण~$\Th{Q}$ में
किसी सूत्र $!B_T$ द्वारा होता है; अर्थात्
$\Th{Q} \Proves !B_T(\num{e}, \num{n}, \num{s})$ तभी और केवल तभी जब
$T(e, n, s)$। यदि
$\Th{Q} \Proves !B_T(\num{e}, \num{n}, \num{s})$, तो
$\Th{Q} \Proves \lexists[y][!B_T(\num{e}, \num{n}, y)]$ भी। यदि ऐसा
कोई~$s$ मौजूद नहीं है, तो प्रत्येक~$s$ के लिए
$\Th{Q} \Proves \lnot !B_T(\num{e}, \num{n}, \num{s})$। किंतु
$\Th{Q}$,~$\omega$-संगत है; अर्थात् यदि प्रत्येक~$n \in \Nat$ के लिए
$\Th{Q} \Proves \lnot !A(\num{n})$, तो
$\Th{Q} \Proves/ \lexists[y][!A(y)]$। हम यह इसलिए जानते हैं क्योंकि
$\Th{Q}$ के अभिगृहीत मानक प्रतिरूप~$\Struct{N}$ में सत्य हैं। अतः
$\Th{Q} \Proves/ \lexists[y][!B_T(\num{e}, \num{n}, y)]$। दूसरे शब्दों
में, $\Th{Q} \Proves \lexists[y][!B_T(\num{e}, \num{n}, y)]$ तभी और
केवल तभी जब कोई $s$ ऐसा हो कि $T(e, n, s)$, अर्थात् तभी और केवल तभी जब
$\cfind{e}(n) \fdefined$। $e$ और~$n$ से हम
$\Gn{\lexists[y][!B_T(\num{e}, \num{n}, y)]}$ संगणित कर सकते हैं; ऐसा
करने वाले आदिम पुनरावर्ती फलन को $g(e, n)$ मानें। अतः
\[
h(e, n) =
\begin{cases}
1 & \text{यदि $\Prov[\Th{Q}](g(e, n))$}\\
0 & \text{अन्यथा}.
\end{cases}
\]
यदि $\Prov[\Th{Q}]$ पुनरावर्ती होता, तो इससे~$h$ भी पुनरावर्ती होता। किंतु
\olref[cmp][rec][hlt]{thm:halting-problem} के अनुसार~$h$ पुनरावर्ती नहीं
है, इसलिए $\Prov[\Th{Q}]$ भी पुनरावर्ती नहीं हो सकता।
\end{proof}

\begin{cor}
प्रथम-कोटि तर्कशास्त्र अनिर्णेय है।
\end{cor}

\begin{proof}
यदि प्रथम-कोटि तर्कशास्त्र निर्णेय होता, तो~$\Th{Q}$ में सिद्धता भी निर्णेय
होती, क्योंकि $\Th{Q} \Proves !A$ तभी और केवल तभी जब
$\Proves !T \lif !A$, जहाँ $!T$,~$\Th{Q}$ के अभिगृहीतों का संयोजन है।
\end{proof}

\end{document}
